\scanpage{47}
$M$ motif sur $k$ \inserted{alg. clos}, $k$ isom. à un sous-corps de $\mathbf C$.

Pour toute place $v$ de $k$, on a $M_v^{\mathbf Z}$ qui est un sous-réseau de $\prod_{\ell\text{ fini}}M(\ell)$. On veut établir les relations entre les $M_v^{\mathbf Z}$ pour diverses valeurs $v_1\ldots v_n$ de $v$. Choisissons un bon foncteur de \unsure{référence} (p.ex. défini par $v_0$)
\[
 \xi:M\rightsquigarrow M_{\unsure{\varphi_0|\xi_0}}\qquad(\unsure{\varphi_0|\xi_0}\in T(\mathbf Q))
\]
avec des isom \unsure{fonctoriels}
\[
 \xi_{0\deleted{\mathbf Q_\ell}}(\unsure{M})_{\mathbf Q_\ell}=M_\ell\otimes_{\mathbf Z_\ell}\mathbf Q_\ell.
\]
Alors $v_1,\ldots,v_n$ définissent des fibrés principaux homogènes
\[
 P_i=\underline{\operatorname{Isom}}(\xi,\xi_{v_i})
\]
sous $G^{\unsure{t|*}}_{M\unsure{\varphi_0|\xi_0}}$ (sur $\mathbf Q$) et on a un point dans
\[
 (\prod P_i)(\widehat{\mathbf Q}),\qquad\text{où }\widehat{\mathbf Q}=\prod_{\ell\text{ fini}}\mathbf Q_\ell,
\]
d'où un pt de $[(\prod P_i)/G](\widehat{\mathbf Q})$.
\margin{$G$ opère diagonalement}
\figure{\begin{tikzpicture}[scale=.5]
\draw (5.1,2.3)--(5.12,2.12) arc (0:-160:.035);
\draw[->] (4.5,2) .. controls (3.8,1.4) and (1.1,.9) .. (0,0);
\end{tikzpicture}}
On peut aussi considérer que l'on a un point canonique dans
\deleted{\[
 \mathrm{Gl}(\unsure{n|m})\times\cdots\times\mathrm{Gl}(\unsure{m|n})
\]}
\deleted{$\mathrm{Gl}$}
\deleted{\[
 \backslash[\mathrm{Gl}(m)\times\cdots\times\mathrm{Gl}(m)](\widehat{\mathbf Q})/\mathrm{Gl}(m)(\widehat{\mathbf Q})
\]}
\deleted{\[
 (\mathrm{Gl}(m)\times\cdots\times\mathrm{Gl}(m))/\mathrm{Gl}(m)
\]}
\[
 \mathrm{Gl}(m,\mathbf Z)\times\cdots\mathrm{Gl}(m,\mathbf Z)\backslash
 \mathrm{Gl}(m,\widehat{\mathbf Q})\times\cdots\times\mathrm{Gl}(m,\widehat{\mathbf Q})/\mathrm{Gl}(m,\widehat{\mathbf Q})
 \qquad m=\operatorname{rang}M
\]
Les considérations qui précèdent montrent que le degré de transcendance du pt en question
\mspage{46}
